% ============================================================================
% What Does MLX 4-bit Cost? A Controlled Audit of Quantization for Code
% Generation on Apple Silicon
%
% Converted from full-draft-v1.md. Every figure traces to an EvalPlus
% per-problem results file; the archived matrices are measured-results-v0313.json
% (unified mlx-lm 0.31.3) and measured-results.json (the earlier 0.29.1 run,
% retained as a cross-version replication).
%
% Build: tectonic -X compile main.tex
% ============================================================================

\documentclass[11pt]{article}

\usepackage[margin=1.25in, top=1in, bottom=1.25in]{geometry}
\usepackage[utf8]{inputenc}
\usepackage[T1]{fontenc}
\usepackage{amsmath}
\usepackage{amssymb}
\usepackage{booktabs}
\usepackage{graphicx}
\usepackage{microtype}
\usepackage{xcolor}
\usepackage{url}
\usepackage[colorlinks=true, linkcolor=blue!70!black,
            citecolor=blue!70!black, urlcolor=blue!70!black,
            pdftitle={What Does MLX 4-bit Cost? A Controlled Audit of Quantization for Code Generation on Apple Silicon},
            pdfauthor={J. Melton, American Code Labs},
            pdfkeywords={quantization, code generation, HumanEval, MBPP, EvalPlus, MLX, Apple Silicon, AWQ, reproducibility},
            pdfsubject={Preprint}]{hyperref}
\usepackage[numbers, sort&compress]{natbib}
\usepackage[htt]{hyphenat}   % allow breaks inside \texttt identifiers
\emergencystretch=2em

\title{What Does MLX 4-bit Cost?\\
A Controlled Audit of Quantization for\\
Code Generation on Apple Silicon}

\author{J.~Melton \\
  American Code Labs \\
  \texttt{jmelton@americancode.org}}

\date{August 2026 \\ \smallskip \small Preprint.}

\begin{document}
\maketitle

% ── Abstract ────────────────────────────────────────────────────────────────
\begin{abstract}
Developers running code models locally on Apple Silicon overwhelmingly use 4-bit
weights from the \texttt{mlx-community} repository, yet no published figure states what
that quantization costs on a code benchmark. Apple reports a single bf16-to-4-bit
accuracy comparison, covering two models on one non-code benchmark; the conversion
model cards carry no accuracy numbers at all. We measure the gap directly.

We evaluate five instruction-tuned code models from 1.5B to 8B parameters on
HumanEval+, MBPP+, and a 26-task suite drawn from our own research codebase, comparing
bf16 against 4-bit weights quantized from those same weights---one harness, one
machine, one variable. Pooled across all five models and both benchmarks, 4-bit
round-to-nearest loses to bf16 on 220 problems and wins on 125 (McNemar
$p < 10^{-6}$), a mean cost of roughly 1 to 6 points that \textbf{shrinks as models
grow} ($+0.52$ points per billion parameters): Phi-4-mini at 3.8B loses 6.5 points
where Qwen2.5-Coder-7B loses 1.2.

Adding an activation-aware (AWQ) arm does not recover the loss. AWQ remains well below
bf16 ($p = 8 \times 10^{-4}$) and is statistically indistinguishable from the
uncalibrated default (133--116, $p = 0.31$), while costing up to 1h51m of quantization
time for identical inference throughput. At these bit widths and scales, \textbf{the
loss is a property of the bit width rather than of the algorithm that produces it.}

We also show that the intuitive way to estimate this cost---differencing a locally
measured 4-bit score against a vendor's published bf16 figure---is unsound: on
Qwen2.5-Coder-1.5B it reports $-9.1$ points where the controlled comparison gives
$-5.5$, because the harness-mismatch term ($-3.6$) is the same order as the effect and
does not even have a stable sign across benchmarks.
\end{abstract}

\paragraph{Completeness of the data.} All 5 of 5 bf16-vs-4-bit pairs are measured, and
4 of 5 carry a calibrated AWQ arm (Phi-4-mini excluded---\texttt{mlx\_lm.quant.awq}
does not support its architecture). Every number below is measured and traceable to an
EvalPlus per-problem results file; the cited figures are archived in
\texttt{measured-results-v0313.json} (unified mlx-lm 0.31.3 runtime) and
\texttt{measured-results.json} (the earlier 0.29.1 matrix, retained as a cross-version
replication). \textbf{No cell is estimated, interpolated, or carried over from a
different configuration.} Where a measurement could not be made---Phi-4-mini's
calibrated arm, and peak memory throughout---the text says so and states why, rather
than substituting a plausible value.

% ── 1. Introduction ─────────────────────────────────────────────────────────
\section{Introduction}
\label{sec:intro}

A developer who wants a code model running locally on a Mac has, in practice, one path:
download a 4-bit conversion from the \texttt{mlx-community} repository and load it with
\texttt{mlx-lm}. The weights are small enough to fit in unified memory, they load in
seconds, and they generate roughly twice as fast as the original bf16 checkpoint. This
is the artifact people actually run.

What it costs them in accuracy is not published anywhere we could find.

The vendor model cards for these conversions report no accuracy numbers; they are
auto-generated conversion artifacts. Apple's own documentation contains exactly one
bf16-to-quantized accuracy comparison---two models, on MMLU Pro, with no code
benchmark~\citep{applebenchmarks}. The one MLX project that does evaluate HumanEval
compares its output against another 4-bit configuration rather than against full
precision, which measures a recipe difference rather than the cost of quantizing at
all. The quantization literature offers a prior of roughly 0.5 to 2 points lost at 4
bits, but every study behind that figure used a \emph{calibrated} method such as
GPTQ~\citep{frantar2023gptq} or AWQ~\citep{lin2024awq}, while MLX's default is
uncalibrated round-to-nearest that additionally quantizes token embeddings and the
output head---a configuration those papers did not measure.

\subsection{Why the number is missing}

The gap is not simply an oversight; it is structurally awkward to close, for three
reasons we encountered rather than anticipated.

First, \textbf{the standard harness does not run on the relevant operating system.}
EvalPlus~\citep{liu2023evalplus} caps executed code via \texttt{RLIMIT\_AS}, which
Darwin rejects, so every problem errors out with a resource-limit exception. Every MLX
user is on macOS. A one-line environment variable works around it
(\S\ref{sec:harness}), but the default experience of trying to evaluate MLX weights
with the field's standard tool is that nothing runs.

Second, \textbf{the obvious shortcut is wrong.} A practitioner who measures their 4-bit
model locally and subtracts the model card's published bf16 score is not measuring
quantization; they are measuring quantization plus every difference between their
harness and the vendor's. On Qwen2.5-Coder-1.5B that shortcut reports a 9.1-point loss
where the controlled comparison gives 5.5. The harness term is 3.6 points---the same
order as the effect---and its sign is not stable: our harness scores below the
published figure on HumanEval and above it on MBPP, so it cannot be corrected with a
constant offset (\S\ref{sec:vendorsub}).

Third, \textbf{the measurement is easy to get wrong quietly.} Building the harness for
this study surfaced five independent defects, each of which scored \emph{correct} model
answers as failures, and none of which announced itself. The largest was a single
undocumented protocol choice---whether the MBPP prompt shows the test assertion, and
therefore conveys the function name the tests require---worth 59 points of pass@1. This
is not incidental to the paper; it is why we abandoned our own harness for EvalPlus,
and it corroborates independent reports that only 12 of 35 published code-model results
could be reproduced (\S\ref{sec:harnessfragility}).

\subsection{What we do}

We evaluate five instruction-tuned code models spanning 1.5B to 8B parameters on a
single 32~GB Apple Silicon host. Each model is measured in bf16 and in 4-bit weights
produced by \texttt{mlx\_lm.convert} \textbf{from those same weights}, so the two arms
differ in exactly one variable. Community uploads are deliberately not used as the
4-bit arm: their provenance cannot be verified, and we show in \S\ref{sec:uploads} that
they are not equivalent to a local conversion. Four of the five models additionally
receive a calibrated AWQ arm; the fifth cannot, because the tool does not support its
architecture---and it is, awkwardly, the model with the largest quantization loss in
the study.

Decoding is greedy with a single sample per problem. That choice requires determinism,
which we verified rather than assumed: 40 of 40 completions were byte-identical across
separate processes (\S\ref{sec:determinism}). This localizes a prior
result~\citep{ouyang2025nondeterminism}---that temperature-zero decoding is
non-deterministic in 47.6\% of HumanEval tasks---as an artifact of API serving rather
than a property of greedy decoding.

Alongside the public benchmarks we evaluate a 26-task suite derived from functions in
our own research codebase. Public code benchmarks are documented as overwhelmingly
easy---one audit finds 84.8\% of HumanEval and 89.6\% of MBPP problems fall in its
lowest difficulty tier~\citep{yadav2024pythonsaga}---and they are demonstrably present
in pretraining corpora. The domain suite is uncontaminated by construction and
discriminates far more sharply across this model range: it spreads the five models by a
factor of 5.5 where HumanEval spreads them by 1.3.

\subsection{What we find}

The cost of 4-bit is real, modest, and larger for smaller models. Pooled across all five
models and both benchmarks, round-to-nearest loses 220 problems to bf16 and wins 125
($p < 10^{-6}$). The per-model mean ranges from $-1.2$ to $-6.5$ points and regresses on
parameter count at $+0.52$ points per billion---meaning the models most often chosen
\emph{because} they are small are the ones quantization hurts most.

Calibration does not fix this. AWQ stays well below bf16 ($p = 8 \times 10^{-4}$) and is
indistinguishable from the uncalibrated default in a head-to-head that has remained at
$p \approx 0.31$ across three successive additions of data. Its cost is paid entirely at
build time---up to 1h51m of quantization for an 8B model---and inference throughput is
identical. The useful conclusion for a practitioner is not ``choose a better
quantizer''; it is that at these bit widths and scales the loss belongs to the bit width
itself.

% ── 2. Related Work ─────────────────────────────────────────────────────────
\section{Related Work}
\label{sec:related}

\subsection{What HumanEval and MBPP actually measure}

The benchmarks this study uses are known to be inadequate in two distinct ways, and
both shape how our results should be read.

\paragraph{Test adequacy.} EvalPlus~\citep{liu2023evalplus} augments HumanEval with
80$\times$ more tests and MBPP with 35$\times$, and finds pass@k falls by 19.3--28.9\%
across 26 models once the additional tests are applied. It also found 18
defects---roughly 11\% of problems---in HumanEval's own reference solutions. We report
both base and plus variants throughout for this reason, and use EvalPlus itself as the
harness rather than a reimplementation (\S\ref{sec:harness}).

\paragraph{Contamination.} Both benchmarks are demonstrably present in common
pretraining corpora: 12.2\% of HumanEval appears verbatim in The Pile and 18.9\% in The
Stack~\citep{matton2024leakage}, with semantic overlap measured at 50.8\% for MBPP and
63.4\% for HumanEval against Stack samples~\citep{riddell2024contamination}. Standard
$n$-gram decontamination has low recall and is defeated by
paraphrase~\citep{yang2023rephrased}. A leak-free replacement benchmark scores models up
to 43\% lower than HumanEval.

This bounds what an absolute score means---but it is \textbf{largely neutralized by our
design}. Both arms of every comparison are the same model on the same problems, so
whatever contamination inflates the bf16 score inflates the 4-bit score too, and cancels
in the paired difference. Contamination is a threat to the \emph{level} of our numbers,
not to the \emph{deltas} that carry our claims.

\paragraph{Difficulty.} PythonSaga's taxonomy audit~\citep{yadav2024pythonsaga} places
84.8\% of HumanEval and 89.6\% of MBPP problems in its lowest difficulty tier, with MBPP
containing none in its highest. Harder or contamination-controlled alternatives
consistently score models far lower: MHPP takes GPT-4o from 91.0\% to
51.1\%~\citep{dai2025mhpp}, EvoEval reports a 39.4\% average drop across 51
models~\citep{xia2024evoeval}, and LiveCodeBench~\citep{jain2025livecodebench} shows the
gap is concentrated in exactly the population we study---fine-tuned open models under
${\sim}7$B, where one model scores 59.8\% on HumanEval+ but 26.3\% on
LiveCodeBench-Easy. This motivates our domain suite (\S\ref{sec:domainsuite}) as a
discriminating instrument rather than a replacement.

\subsection{Quantization and code generation}

The prior for 4-bit weight-only quantization is a loss of roughly 0.5 to 2 pass@1 points
in the 1B--34B range. \citet{giagnorio2025ptq} measure CodeLlama-7B at
$29.8 \to 29.1$ on MultiPL-E Python going from FP16 to 4-bit with AQLM, with the sharp
degradation appearing only at 3 bits and below. \citet{kurtic2025bf16}, across more than
500,000 evaluations, report W4A16 recovering 98.9\% of BF16 HumanEval performance.
\citet{fang2025smaller}, whose model range most closely matches ours, find 4-bit costs
1.19\% of clean performance on average and conclude that smaller models tolerate
quantization better than expected.

Not all results agree. \citet{shi2025systematic} report far larger degradation and argue
coding and math are disproportionately fragile---though their catastrophic cases are
dominated by activation quantization rather than the weight-only setting we study.
Separately, scaling work finds low-bit degradation grows with training tokens and
shrinks with model size~\citep{ouyang2024lowbit}, which places modern heavily-trained
1--8B models in the least favourable quadrant. \textbf{Our size trend
(\S\ref{sec:paired}) is consistent with that prediction}; our absolute magnitudes sit at
or slightly above the calibrated-method prior, which is what one would expect given that
MLX's default is uncalibrated (\S\ref{sec:quantprovenance}).

Critically, \textbf{every study above measures a calibrated quantizer.} None measures
naive round-to-nearest with embeddings and the output head also quantized, which is what
\texttt{mlx\_lm.convert -{}-q-bits 4} produces and what the artifacts on
\texttt{mlx-community} are.

\subsection{Apple Silicon evaluation}

Published Apple Silicon LLM benchmarking is overwhelmingly about speed and memory rather
than accuracy. Apple's own \texttt{mlx-lm} benchmark
documentation~\citep{applebenchmarks} contains the only first-party bf16-to-quantized
accuracy comparison we located: two models, on MMLU Pro, with a bf16$\to$4-bit delta of
$-3.33$ points for a 4B model. Crowdsourced Apple Silicon performance tables and recent
throughput studies report tokens/sec and memory across many models and report no
accuracy figures at all. The \texttt{mlx-community} conversion cards are auto-generated
and carry none either.

\textbf{We could find no published HumanEval or MBPP result for stock
\texttt{mlx-community} 4-bit weights against a bf16 baseline}, from Apple, from
mlx-community, or from a third party. That intersection---MLX's default quantizer
$\times$ code benchmarks $\times$ the 1--8B range---is the cell this paper fills.

\subsection{Reproducibility of code-model results}

Our harness findings (\S\ref{sec:harnessfragility}) sit in a small but consistent
literature. \citet{szalontai2025repro} attempted to reproduce 35 published results on a
HumanEval variant and succeeded on 12, attributing failures to benchmark-variant
confusion, temperature misconfiguration, and precision that papers never state.
\citet{ouyang2025nondeterminism} report that 47.56\% of HumanEval tasks produce
non-identical outputs across repeated requests even at temperature zero---a result we
localize in \S\ref{sec:determinism} as an artifact of API serving rather than of greedy
decoding. FormatSpread~\citep{sclar2024formatspread} demonstrates up to 76 accuracy
points of spread from prompt formatting alone, albeit on classification rather than
code.

What appears to be missing is a code-native account of how much the \emph{harness} moves
pass@1. \S\ref{sec:harnessfragility} contributes five measured instances, the largest
worth 59 points.

% ── 3. Methods ──────────────────────────────────────────────────────────────
\section{Methods}
\label{sec:methods}

\subsection{Design}

One variable. For each model we evaluate two arms---bf16 and 4-bit---through an
identical harness, on the same host, with identical decoding. Vendor published numbers
are used \textbf{only} as a sanity check that our harness reproduces the standard
protocol, never as a term in any difference.

\subsection{Harness}
\label{sec:harness}

EvalPlus 0.3.1~\citep{evalplusrepo} (HumanEval+ 164 tasks, MBPP+ 378 tasks; the MBPP+
task count was revised from 399 in an earlier dataset version), rather than a bespoke
harness.
This is a deliberate reversal: we began with a custom implementation and abandoned it
after finding five independent defects (\S\ref{sec:harnessfragility}).

\textbf{EvalPlus does not run on macOS as shipped.} \texttt{reliability\_guard} caps
executed code via \texttt{RLIMIT\_AS}, which Darwin rejects, and every problem fails with
\texttt{ValueError: current limit exceeds maximum limit}. The documented escape hatch is
\texttt{EVALPLUS\_MAX\_MEMORY\_BYTES=-1}. We note this because every MLX user is on
macOS, and it plausibly contributes to the absence of published MLX code-benchmark
numbers.

Generation runs under system Python 3.9 (where \texttt{mlx-lm} is installed) and scoring
under Python 3.14 (where EvalPlus installs cleanly); EvalPlus decouples the two by
design.

\subsection{Quantization and provenance}
\label{sec:quantprovenance}

4-bit arms are produced \textbf{locally} by
\texttt{mlx\_lm.convert -{}-quantize -{}-q-bits 4} from the same weights the bf16 arm
evaluates, rather than downloaded from \texttt{mlx-community}. A community upload cannot
be verified as to base revision or conversion settings, which would confound the single
variable under test. \S\ref{sec:uploads} shows this is not a pedantic concern.

MLX's default quantizer is affine round-to-nearest~\citep{applequantize}: per group of
64, scale and bias from min/max, no calibration data, no error compensation. Both
\texttt{nn.Linear} and \texttt{nn.Embedding} define \texttt{to\_quantized}, so
\textbf{token embeddings and \texttt{lm\_head} are quantized to 4 bits as
well}---where GGUF Q4\_K\_M keeps them at Q6\_K and GPTQ/AWQ pipelines conventionally
leave them in fp16. Nominal ``4-bit'' is therefore not one thing.

\subsection{Decoding and the determinism gate}
\label{sec:determinism}

Greedy, \texttt{temp=0.0}, $n=1$, max 768 new tokens. $n=1$ requires determinism, which
we verified rather than assumed: 20 HumanEval problems generated twice within a process
and twice across separate processes produced \textbf{40/40 byte-identical completions}.

This matters because \citet{ouyang2025nondeterminism} report that 47.6\% of HumanEval
tasks yield non-identical outputs at temperature 0. Our result localizes that finding: it
is an artifact of API serving (batching, load balancing, kernel selection), not a
property of greedy decoding. Local MLX greedy is bit-reproducible. Throughput varied
193--208 tok/s across those identical runs---speed is noisy, tokens are not.

\subsection{Statistics}

Arms are evaluated on identical problems, so observations are paired. We use McNemar's
exact test on discordant pairs, not a difference of proportions. We report $b$ (bf16
solved, 4-bit failed) and $c$ (the reverse) throughout.

\subsection{Domain task suite}
\label{sec:domainsuite}

Public benchmarks are saturated and skewed easy---PythonSaga finds 84.8\% of HumanEval
and 89.6\% of MBPP problems are ``Easy'', with MBPP containing no ``Hard'' problems at
all~\citep{yadav2024pythonsaga}. We therefore add 26 tasks derived from functions in our
own research codebase (SAE training, circuit tracing, activation streaming, statistical
nulls): 8 EASY, 8 MEDIUM, 10 HARD, 347 assertions.

Each task was validated in both directions: it must pass with a known-correct reference
solution, \textbf{and} reject 21 hand-written plausible-but-wrong implementations. This
two-sided validation caught one task whose stated trap was mathematically vacuous---it
would have scored every model correct---and one that was unanswerable by construction.

\subsection{Host}

Apple Silicon, 32~GB unified memory. This ceiling excludes DeepSeek-Coder-V2-Lite from
the bf16 arm (31.4~GB weights); it is reported 4-bit-only and contributes to no delta.

% ── 4. Results: quantization cost ───────────────────────────────────────────
\section{Results: Quantization Cost}
\label{sec:results}

\subsection{Paired comparison}
\label{sec:paired}

All five pairs are measured. $\Delta$ is the mean of the HumanEval and MBPP deltas; $b$
and $c$ are discordant pairs pooled over both benchmarks (bf16-only and 4-bit-only
solves).

\begin{table}[h]
  \centering
  \caption{Paired bf16-vs-4-bit comparison, all five models. $\Delta$ is in percentage
    points. $\dagger$~Qwen-3B's 4-bit arm is an \texttt{mlx-community} upload, not
    locally converted (\S\ref{sec:uploads}); excluding it leaves 175 discordant pairs
    favouring bf16 against 102 favouring 4-bit, so neither the sign nor the significance
    of the pooled result changes.}
  \label{tab:paired}
  \smallskip
  \scriptsize
  \setlength{\tabcolsep}{4pt}
  \begin{tabular}{lrrrrrrrrr}
    \toprule
    Model & Params & HE bf16 & HE 4bit & MBPP bf16 & MBPP 4bit & mean $\Delta$ & $b$ & $c$ & $p$ \\
    \midrule
    Qwen2.5-Coder-1.5B      & 1.5B & 0.671 & 0.616 & 0.704 & 0.680 & $-3.9$ & 35 & 17 & \textbf{0.018} \\
    Qwen2.5-Coder-3B$^\dagger$ & 3.0B & 0.841 & 0.811 & 0.772 & 0.728 & $-3.8$ & 45 & 23 & \textbf{0.010} \\
    Phi-4-mini              & 3.8B & 0.665 & 0.616 & 0.646 & 0.563 & $\mathbf{-6.5}$ & 91 & 52 & \textbf{0.0014} \\
    Qwen2.5-Coder-7B        & 7.0B & 0.896 & 0.878 & 0.849 & 0.844 & $-1.2$ & 16 & 11 & 0.442 \\
    Granite-8B-Code         & 8.0B & 0.591 & 0.579 & 0.672 & 0.648 & $-1.8$ & 33 & 22 & 0.177 \\
    \midrule
    \textbf{All five pooled} & & & & & & & \textbf{220} & \textbf{125} & $\mathbf{< 10^{-6}}$ \\
    \bottomrule
  \end{tabular}
\end{table}

\paragraph{Two results, not one.}
\begin{enumerate}
  \item \textbf{The cost is real.} 220 discordant pairs favour bf16 against 125
    favouring 4-bit, $p < 10^{-6}$. Individually only three of five models reach
    significance, which is why single-model studies of this effect are
    underpowered---an important observation in its own right.
  \item \textbf{The cost shrinks with model size.} Regressing mean $\Delta$ on parameter
    count gives \textbf{$+0.52$ points per billion}: Phi-4-mini (3.8B) loses 6.5 points
    while Qwen-7B loses 1.2. This is consistent with the quantization literature's
    finding that larger models are more resilient~\citep{ouyang2024lowbit}, and it is the
    practically actionable result---the models most often chosen \emph{because} they are
    small are the ones 4-bit hurts most.
\end{enumerate}

\textbf{Phi-4-mini is the outlier in both directions}---the largest quantization cost in
the matrix ($-6.5$) and the largest gap to its published bf16 figure ($-7.9$,
\S\ref{sec:baselines}). We do not have an explanation. Two candidates are consistent
with the evidence but untested here: its published figure comes from an internal
pipeline using 3-shot MBPP where we use 0-shot~\citep{microsoft2025phi4mini}, which would
inflate the reported baseline without affecting our paired delta; and MLX quantizes token
embeddings and the output head (\S\ref{sec:quantprovenance}), so a model whose embedding
matrix is unusually large relative to its body would be disproportionately affected.
Distinguishing these would require an arm with embeddings held at higher precision, which
\texttt{mlx\_lm.convert} supports only through a non-default mixed recipe. We flag it
rather than resolve it.

\subsection{bf16 baselines against published figures}
\label{sec:baselines}

Sanity check only---never differenced against a 4-bit score. Table~\ref{tab:baselines}
gives our bf16 HumanEval scores beside the published figures.

\begin{table}[h]
  \centering
  \caption{Our bf16 HumanEval scores against published figures.
    $\ddagger$~IBM publishes HumanEvalSynthesize (57.9 Python) rather than vanilla
    HumanEval~\citep{ibm2024granite}; our 59.1 is near it but not the same benchmark.}
  \label{tab:baselines}
  \smallskip
  \begin{tabular}{lrrr}
    \toprule
    Model & Ours (bf16) & Published & $\Delta$ \\
    \midrule
    Qwen2.5-Coder-1.5B & 0.671 & 0.707 & $-3.6$ \\
    Qwen2.5-Coder-3B   & 0.841 & 0.841 & \textbf{0.0} \\
    Qwen2.5-Coder-7B   & 0.896 & 0.884 & $+1.2$ \\
    Phi-4-mini         & 0.665 & 0.744 & $-7.9$ \\
    Granite-8B-Code    & 0.591 & n/a$^\ddagger$ & --- \\
    \bottomrule
  \end{tabular}
\end{table}

Agreement is close across the Qwen family~\citep{qwen2024coder} and exact at 3B, which is
the strongest available evidence that the harness reproduces the vendor protocol.
Phi-4-mini's larger gap is consistent with Microsoft evaluating through an internal
pipeline whose prompt format and shot count are unspecified.

\subsection{Throughput}
\label{sec:throughput}

Table~\ref{tab:throughput} reports generation throughput across all three arms.

\begin{table}[h]
  \centering
  \caption{Generation throughput (tok/s) during the HumanEval sweep, all arms, under the
    unified mlx-lm 0.31.3 runtime. Speedup is 4-bit RTN over bf16.}
  \label{tab:throughput}
  \smallskip
  \begin{tabular}{lrrrr}
    \toprule
    Model & bf16 & 4-bit RTN & 4-bit AWQ & speedup \\
    \midrule
    Qwen2.5-Coder-1.5B & 90.6 & 164.8 & 164.8 & 1.82$\times$ \\
    Qwen2.5-Coder-3B   & 48.1 &  97.8 &  97.5 & 2.03$\times$ \\
    Phi-4-mini         & 41.0 &  88.9 & n/a   & 2.17$\times$ \\
    Qwen2.5-Coder-7B   & 22.8 &  57.3 &  57.2 & 2.51$\times$ \\
    Granite-8B-Code    & 19.7 &  48.3 &  48.2 & 2.45$\times$ \\
    \bottomrule
  \end{tabular}
\end{table}

4-bit is 1.8--2.5$\times$ faster, and \textbf{the speedup grows with model size}---the
opposite of the accuracy cost, which shrinks with size (\S\ref{sec:paired}). The two
trends compound in the same direction: larger models gain more throughput from
quantization and lose less accuracy to it.

\textbf{RTN and AWQ are throughput-identical} to within measurement noise (164.8 vs
164.8; 57.3 vs 57.2; 48.3 vs 48.2). Whatever the calibrated quantizer costs is paid
entirely at build time (\S\ref{sec:calibration}), never at inference.

\paragraph{Memory was not measured reliably, and no memory claim is made.} Our collection
path reports 4-bit consuming \emph{more} peak memory than bf16 on Qwen-3B (6.55 vs
6.47~GB) while being sane on Qwen-1.5B (1.65 vs 3.48~GB)---an internal contradiction we
did not resolve. Since the memory reduction is half the practical motivation for 4-bit,
this is a substantive limitation rather than a footnote, and is restated in
\S\ref{sec:threats}. Readers should take the well-established weight-footprint reduction
(${\sim}4\times$ on the quantized tensors) as the basis for memory reasoning, not any
figure from this study.

\subsection{Calibration does not rescue 4-bit}
\label{sec:calibration}

MLX's default quantizer is uncalibrated round-to-nearest (\S\ref{sec:quantprovenance}),
which is weaker than the GPTQ/AWQ-class methods behind the published
${\sim}0.5$--2~pt loss figures. The obvious hypothesis is that the cost we measure is an
artifact of that default rather than a property of 4-bit. We tested it directly by adding
an activation-aware (AWQ) arm~\citep{lin2024awq} at the same bit width and group size.

Table~\ref{tab:threearm} gives the three-arm comparison. All three arms ran under a single mlx-lm 0.31.3 (see \S\ref{sec:runtime} on why the
runtime changed and why it is not a confound). \textbf{Four of five models carry a
calibrated arm; Phi-4-mini cannot ($\ddagger$).}

\begin{table}[h]
  \centering
  \caption{Three-arm comparison. $\Delta$ columns are percentage points against bf16.}
  \label{tab:threearm}
  \smallskip
  \small
  \setlength{\tabcolsep}{5pt}
  \begin{tabular}{llrrrrr}
    \toprule
    Model & Benchmark & bf16 & 4-bit RTN & 4-bit AWQ & RTN $\Delta$ & AWQ $\Delta$ \\
    \midrule
    Qwen2.5-Coder-1.5B & HumanEval & 0.683 & 0.622 & 0.628 & $-6.1$ & $-5.5$ \\
                       & MBPP      & 0.704 & 0.688 & 0.683 & $-1.6$ & $-2.1$ \\
    Qwen2.5-Coder-3B   & HumanEval & 0.854 & 0.829 & 0.835 & $-2.4$ & $-1.8$ \\
                       & MBPP      & 0.775 & 0.712 & 0.738 & $-6.3$ & $-3.7$ \\
    Granite-8B-Code    & HumanEval & 0.598 & 0.585 & 0.549 & $-1.2$ & $-4.9$ \\
                       & MBPP      & 0.677 & 0.648 & \textbf{0.680} & $-2.9$ & $\mathbf{+0.3}$ \\
    Qwen2.5-Coder-7B   & HumanEval & 0.896 & 0.878 & 0.890 & $-1.8$ & $-0.6$ \\
                       & MBPP      & 0.854 & 0.844 & 0.841 & $-1.1$ & $-1.3$ \\
    Phi-4-mini         &           & 0.665 & 0.616 & \textbf{n/a}$^\ddagger$ & $-4.9$ & \\
    \bottomrule
  \end{tabular}
\end{table}

$\ddagger$~\textbf{Phi-4-mini cannot have an AWQ arm.} \texttt{mlx\_lm.quant.awq} raises
\texttt{NotImplementedError: AWQ support for phi3 models NYI}. This is a hard limitation
of the tool, not a transient failure. It matters because Phi-4-mini carries the
\emph{largest} RTN cost in the matrix (\S\ref{sec:paired}), so it is precisely the model
where ``a better quantizer would fix this'' was most plausible---and it is the one model
where we cannot test it. Reported as excluded with cause rather than silently dropped.

Table~\ref{tab:mcnemar} pools McNemar over the 8 cells where all three arms exist (4
models $\times$ 2 benchmarks; Phi-4-mini is absent because it has no AWQ arm). These
counts are therefore a subset of the five-model totals in \S\ref{sec:paired} and are not
directly comparable to them.

\begin{table}[h]
  \centering
  \caption{Pooled McNemar over the 8 three-arm cells (4 models $\times$ 2 benchmarks).}
  \label{tab:mcnemar}
  \smallskip
  \begin{tabular}{lrrr}
    \toprule
    Comparison & wins A & wins B & $p$ \\
    \midrule
    bf16 vs RTN & 134 (bf16) & 70 (RTN) & \textbf{0.000009} \\
    bf16 vs AWQ & 119 (bf16) & 72 (AWQ) & \textbf{0.00083} \\
    \textbf{RTN vs AWQ (head-to-head)} & 133 (AWQ) & 116 (RTN) & \textbf{0.31} \\
    \bottomrule
  \end{tabular}
\end{table}

\textbf{Both quantizers lose to bf16 by a wide statistical margin. Neither beats the
other.} The head-to-head has been stable at $p \approx 0.31$ across three successive
additions of data, and Qwen-7B's own head-to-head is 14--13 ($p = 1.00$)---as close to a
tie as the test can express.

The practical conclusion is therefore not ``use a better quantizer''. It is that at 4
bits and these scales \textbf{the loss is a property of the bit width, not of the
algorithm that gets you there.}

\paragraph{Per-benchmark variance exceeds the method difference.} Granite-8B is the
clearest case: AWQ is 4.9 points \emph{below} bf16 on HumanEval but \emph{matches} it on
MBPP ($+0.3$), while RTN loses on both. On the same model with the same weights, the two
benchmarks disagree about which quantizer is better. This is the section's central
caution. Read one benchmark on one model and the directional conclusion reverses---the
same failure this paper documents in \S\ref{sec:harnessfragility}, arriving from a
different direction. \textbf{A single benchmark on a single model does not support a
directional claim about quantizer quality.}

The head-to-head also bears on \S\ref{sec:uploads}, where an \texttt{mlx-community}
upload outscored a local RTN conversion and appeared to suggest that recipe quality
matters. AWQ versus RTN is a far larger recipe difference, measured on four models, and
it finds no significant effect ($p = 0.31$). Both can be true---the community uploads may
differ from local \texttt{convert} in some way other than calibration---but the simpler
reading is that \S\ref{sec:uploads}'s effect is noise, and it is reported as
non-significant there.

\paragraph{Cost of the calibrated arm.} AWQ quantization scales worse than linearly: 18
minutes at 1.5B, ${\sim}$1h30m at 7B, 1h51m at 8B---roughly 6$\times$ the time for
5.3$\times$ the parameters, as the grid search over scaling factors compounds with layer
count. Inference throughput is unchanged (Granite-8B 48.2 vs 48.3 tok/s; Qwen-7B 57.2 vs
57.3), so the price is paid entirely at build time and buys no runtime benefit either.

\paragraph{Domain suite.} The 26-task suite tracks the same ordering but with a wider
spread, and is the one place AWQ looks consistently worse (Table~\ref{tab:domainarms}).

\begin{table}[h]
  \centering
  \caption{Domain task suite, three arms, tasks solved of 26.}
  \label{tab:domainarms}
  \smallskip
  \begin{tabular}{lrrr}
    \toprule
    Model & bf16 & RTN & AWQ \\
    \midrule
    Qwen-7B    & 12/26 & 11/26 & \textbf{8/26} \\
    Granite-8B &  7/26 &  6/26 & 7/26 \\
    \bottomrule
  \end{tabular}
\end{table}

Qwen-7B's AWQ arm solves four fewer tasks than bf16 while its HumanEval score is within
0.6 points. With only 26 tasks this is a small-$N$ observation, not a result---but it is
the direction worth checking if the suite is expanded, since it hints that aggregate
pass@1 on saturated benchmarks may hide capability loss the harder tasks expose.

\paragraph{One asymmetry to record rather than smooth over.} AWQ sets
\texttt{model.embed\_tokens} to \texttt{group\_size:~32} while RTN uses 64 uniformly. So
``RTN vs AWQ'' is not a pure method comparison---the embedding treatment differs too.
Given \S\ref{sec:quantprovenance} (MLX quantizes embeddings and \texttt{lm\_head} at all,
unlike GGUF Q4\_K\_M and GPTQ/AWQ pipelines elsewhere), this is a plausible contributor
and should not be attributed to calibration alone.

\subsection{Runtime version is not a confound}
\label{sec:runtime}

The calibrated quantizers exist only in mlx-lm $\geq$ 0.30~\citep{applelearnedquants},
while the matrix in \S\ref{sec:paired} was measured under 0.29.1. The original plan was
to produce AWQ weights with a newer mlx-lm in an isolated venv and evaluate them on the
pinned runtime, keeping inference constant. That failed: 0.31.3 writes a per-layer
quantization config that 0.29.1 cannot parse. We rejected writing a compatibility shim,
because misreading a quantization spec would produce an AWQ column that looks correct and
means something else.

Instead we tested the assumption being protected (Table~\ref{tab:runtime}). Same weights,
same harness, only the mlx-lm version differing.

\begin{table}[h]
  \centering
  \caption{Runtime-version control on Qwen-1.5B. 162 of 164 per-problem verdicts are
    identical.}
  \label{tab:runtime}
  \smallskip
  \begin{tabular}{lrrrlr}
    \toprule
    & 0.29.1 & 0.31.3 & $\Delta$ & discordant & $p$ \\
    \midrule
    Qwen-1.5B HumanEval  & 0.671 & 0.683 & $+1.2$ & 0 old-only / 2 new-only & 0.50 \\
    Qwen-1.5B HumanEval+ & 0.622 & 0.634 & $+1.2$ & 0 / 2 & 0.50 \\
    \bottomrule
  \end{tabular}
\end{table}

\textbf{162 of 164 per-problem verdicts are identical.} The runtime shifts 2 problems;
the effect under study accounts for 345 discordant pairs across the matrix. The version
is therefore not a confound at the resolution that matters, and the three-arm results run
under a single 0.31.3.

The 0.29.1 matrix is retained and reported (\S\ref{sec:paired}) rather than discarded---it
is a cross-version replication, which is more than the design originally promised. Every
summary records \texttt{mlx\_lm\_version}, \texttt{python}, and \texttt{runs\_base} so no
score can be silently attributed to the wrong inference stack.

% ── 5. Vendor subtraction ───────────────────────────────────────────────────
\section{Results: The Vendor-Subtraction Confound}
\label{sec:vendorsub}

The intuitive estimate of quantization cost---measure 4-bit locally, subtract the model
card's bf16 figure---is available to any practitioner and is wrong
(Table~\ref{tab:vendorsub}).

\begin{table}[h]
  \centering
  \caption{Two estimates of the same quantity, Qwen-1.5B HumanEval.}
  \label{tab:vendorsub}
  \smallskip
  \begin{tabular}{lr}
    \toprule
    Method & Implied cost \\
    \midrule
    Naive: published bf16 (0.707) $-$ measured 4-bit (0.616)   & $\mathbf{-9.1}$ pts \\
    Controlled: measured bf16 (0.671) $-$ measured 4-bit (0.616) & $\mathbf{-5.5}$ pts \\
    \bottomrule
  \end{tabular}
\end{table}

The naive figure absorbs the entire harness-mismatch term ($-3.6$ pts here), which is of
the same order as the effect. The sign of that term is not even stable: our harness
scores \emph{below} published on HumanEval for Qwen-1.5B but \emph{above} published on
MBPP, ruling out a simple ``our prompt is worse'' explanation and showing the bias cannot
be corrected with a constant.

The same comparison for Qwen2.5-Coder-3B, whose 4-bit arm is provenance-matched in the
unified run: published bf16 0.841, measured bf16 0.854, measured 4-bit 0.829. The naive
estimate gives $-1.2$ and the controlled one $-2.5$---here the shortcut \emph{understates}
the cost, because our harness scores above the published figure on this model. The error
does not have a consistent direction, which is precisely why it cannot be corrected for.

% ── 6. Community uploads ────────────────────────────────────────────────────
\section{Observation: Community Uploads May Not Equal Local Conversions}
\label{sec:uploads}

\textbf{This section reports an observation, not a result.} It is retained because it
motivates a methodological choice (\S\ref{sec:quantprovenance}) and because the
alternative---omitting a measurement that did not reach significance---is the kind of
selective reporting this paper argues against elsewhere.

Early in this study the 4-bit arm was an \texttt{mlx-community} upload rather than a local
conversion. Measuring both for Qwen2.5-Coder-1.5B, evaluated identically
(Table~\ref{tab:uploads}):

\begin{table}[h]
  \centering
  \caption{Community upload against local conversion, Qwen2.5-Coder-1.5B.}
  \label{tab:uploads}
  \smallskip
  \begin{tabular}{lrrrr}
    \toprule
    Benchmark & community upload & self-converted & $\Delta$ & $p$ \\
    \midrule
    HumanEval & 0.646 & 0.616 & $-3.0$ & 0.18 \\
    MBPP      & 0.690 & 0.680 & $-1.1$ & 0.46 \\
    pooled    &       &       &        & \textbf{0.11} \\
    \bottomrule
  \end{tabular}
\end{table}

The upload scores higher on both benchmarks, but neither cell nor the pooled comparison
reaches significance on a single model.

\paragraph{Why we do not claim this.} Two things argue against promoting it. First, it
was never replicated: we switched to local conversion for provenance control
(\S\ref{sec:quantprovenance}) and did not measure a second community upload, so this
rests on one model where $p = 0.11$. Second, it sits in tension with
\S\ref{sec:calibration}. This observation suggests the conversion recipe matters; the AWQ
head-to-head---a far larger recipe difference, on four models---found no significant
effect ($p = 0.31$). The simplest reading consistent with both is that this section's
difference is noise.

\paragraph{What it does justify} is the methodological choice. Whether or not the
difference is real, a community upload's provenance cannot be verified: the base revision
and conversion settings are not recorded on the auto-generated cards. Using one as the
4-bit arm would place an unverifiable artifact on one side of a comparison whose entire
value rests on differing in a single variable. Converting locally removes that risk at no
cost, and we recommend it for any study of this kind.

\paragraph{If someone wants to settle it}, the experiment is cheap: evaluate the community
upload and a local conversion of the same base weights across several models. It is worth
settling, because if uploads really are better then every number in this paper is a lower
bound on what practitioners actually deploy---they run the uploads, not local conversions.

% ── 7. Harness fragility ────────────────────────────────────────────────────
\section{Harness Fragility as a Finding}
\label{sec:harnessfragility}

Five independent defects were found in a custom HumanEval/MBPP harness before it was
abandoned (Table~\ref{tab:defects}). Every one silently scored \textbf{correct} model
answers as failures.

\begin{table}[h]
  \centering
  \caption{Five harness defects, each of which scored correct answers as failures.}
  \label{tab:defects}
  \smallskip
  \small
  \begin{tabular}{p{0.5\textwidth}p{0.36\textwidth}}
    \toprule
    Defect & Measured effect \\
    \midrule
    MBPP prompt never conveyed the required function name, which \texttt{test\_list}
      asserts exactly & MBPP pass@1 \textbf{0.047 $\to$ 0.640} \\
    One system prompt demanded a bare body---correct for HumanEval's stub, fatal for
      MBPP which has none & SyntaxError on every compliant answer \\
    \texttt{test\_imports} never emitted & 13/257 unconditional NameErrors \\
    Unindented body appended to a HumanEval stub & SyntaxError on correct answers \\
    Extractor's non-greedy fence match truncated any answer containing a code-fence
      delimiter & one task unanswerable by construction \\
    \bottomrule
  \end{tabular}
\end{table}

The first defect is the headline: a single undocumented protocol choice moved pass@1 by
59 points. EvalPlus's own MBPP+ prompts include the assertion, so this was a deviation
from the accepted protocol rather than a hard benchmark. This corroborates
\citet{szalontai2025repro}, who reproduced only 12 of 35 published results and attributed
failures to benchmark-variant confusion, temperature misconfiguration, and unspecified
precision.

The practical rule this suggests: \textbf{a benchmark number is not evidence until the
harness has been shown to score a known-correct reference answer as passing.} Our domain
suite enforces this in both directions (\S\ref{sec:domainsuite}).

These belong in the body rather than an appendix. The reproducibility literature
establishes that harness configuration moves scores---\citet{szalontai2025repro}
reproduced 12 of 35 published results, FormatSpread~\citep{sclar2024formatspread} measured
76 points of spread from formatting---but those are, respectively, a variant benchmark and
a classification task. We could find no code-native measurement of how much the harness
itself moves pass@1. The first defect is exactly that measurement: 59 points, from a
protocol choice that papers do not document. A result of that magnitude, in the same units
as every score this paper reports, is not supporting material.

% ── 8. Domain suite ─────────────────────────────────────────────────────────
\section{Results: Domain Task Suite}
\label{sec:domainresults}

Table~\ref{tab:domain} sets the domain suite beside HumanEval for the bf16 arms.

\begin{table}[h]
  \centering
  \caption{Domain task suite against HumanEval, bf16 arms. DeepSeek-V2-Lite is
    4-bit-only (its bf16 weights exceed the host's 32~GB) and contributes to no delta.}
  \label{tab:domain}
  \smallskip
  \begin{tabular}{lrr}
    \toprule
    Model & Domain suite & HumanEval \\
    \midrule
    Qwen2.5-Coder-1.5B        & 2/26  & 0.671 \\
    Qwen2.5-Coder-3B          & 5/26  & 0.841 \\
    Granite-8B-Code           & 6/26  & 0.591 \\
    Phi-4-mini                & 7/26  & 0.665 \\
    DeepSeek-V2-Lite (4-bit)  & 7/26  & 0.762 \\
    Qwen2.5-Coder-7B          & \textbf{11/26} & 0.896 \\
    \bottomrule
  \end{tabular}
\end{table}

The suite spreads this model range 5.5$\times$ (2 $\to$ 11) where HumanEval spreads it
1.3$\times$ (0.671 $\to$ 0.896). The best model solves 42\% of tasks; the smallest solves
8\%. This is consistent with the PythonSaga finding that public benchmarks are dominated
by easy problems~\citep{yadav2024pythonsaga}, and supports using an uncontaminated domain
suite as the discriminating instrument.

A per-tier and per-trap-family breakdown across arms is possible from the recorded
per-task results, but at 26 tasks split three ways the cells are too small to support
inference---a single task moves a tier by 10--12 percentage points. We report the
aggregate only, and note that expanding the suite is the prerequisite for any
finer-grained claim (\S\ref{sec:next}).

% ── 9. Discussion ───────────────────────────────────────────────────────────
\section{Discussion}
\label{sec:discussion}

\subsection{What this means if you are choosing weights}

For most local coding work, \textbf{4-bit is the right default, and the reason is the
exchange rate rather than the loss being negligible.} The loss is real and statistically
unambiguous ($p < 10^{-6}$ pooled), but it buys roughly a doubling of throughput and
roughly half the memory. On a 32~GB machine that trade is often what makes a model usable
at all: Qwen2.5-Coder-7B in bf16 leaves little headroom, while its 4-bit conversion runs
comfortably alongside an editor and a browser.

The more actionable finding is \textbf{where the cost falls.} It shrinks as models
grow---$+0.52$ points per billion parameters---so the penalty is largest exactly where
users are most likely to accept it. Phi-4-mini at 3.8B loses 6.5 points;
Qwen2.5-Coder-7B loses 1.2. A practitioner choosing a small model \emph{because} it is
small is paying the steepest quantization tax, and would often be better served by a
larger model at 4 bits than a smaller one at higher precision---a conclusion consistent
with prior work finding that at equal memory, bigger-model-lower-bit beats
smaller-model-higher-bit~\citep{jin2024comprehensive}.

\subsection{``Use a better quantizer'' is not the fix}

The intuitive response to a quantization loss is to reach for a better quantization
method. Our data does not support that move. AWQ---activation-aware, calibrated, the
class of method behind the published ${\sim}$1-point figures---remains well below bf16
($p = 8 \times 10^{-4}$) and is statistically indistinguishable from the naive default
(133--116, $p = 0.31$). That head-to-head held at $p \approx 0.31$ across three
successive additions of data, so it is stable rather than merely underpowered, and
Qwen-7B's own head-to-head is 14--13.

The cost side makes this decisive rather than merely inconclusive. AWQ quantization took
up to 1h51m for an 8B model, scaling worse than linearly with parameter count, and
produced \textbf{identical inference throughput} (57.2 vs 57.3 tok/s). A practitioner
would be spending hours of compute for a benefit this study cannot detect.

We state the conclusion as: at 4 bits and these scales, \textbf{the loss is a property of
the bit width, not of the algorithm that produces it.} We are careful not to overreach.
This is one calibrated method on four models; GPTQ and DWQ are unmeasured, and AWQ and
RTN differ in embedding group size (32 vs 64) as well as in calibration, so the
head-to-head is not a pure method contrast.

\subsection{Do not compare your numbers to a model card}

The single most transferable result here is negative. Differencing a locally measured
4-bit score against a vendor's published bf16 figure is a natural thing to do and it does
not work: on Qwen2.5-Coder-1.5B it reports $-9.1$ points where the controlled comparison
gives $-5.5$. The harness-mismatch term is 3.6 points---the same order as the
effect---and its sign is not stable across benchmarks, so it cannot be absorbed into a
correction factor.

The corollary is that a bf16 baseline is not optional. It is the majority of the
experimental cost in this study, and it is the only part that makes the 4-bit number mean
anything.

\subsection{The measurement is the hard part}

Five of the defects we found (\S\ref{sec:harnessfragility}) scored \emph{correct} answers
as failures, and none announced itself. The largest---whether the MBPP prompt conveys the
function name its tests assert---was worth 59 points of pass@1, comfortably larger than
every effect this paper reports. That we found it only because a 4.7\% score was too
implausible to accept is not reassuring; a defect of half that size would have produced a
believable number and been published.

This shaped our practice in ways we would recommend generally. Every task in the domain
suite must pass with a known-correct reference \textbf{and} reject 21 hand-written wrong
implementations, a two-sided check that caught one task whose trap was mathematically
vacuous (it would have scored every model correct) and one that was unanswerable by
construction. Scores are read from the harness's own per-problem verdict files rather than
from any summary a runner wrote. And a benchmark sweep that returns 0/N is treated as a
harness fault rather than a result.

The same discipline applies to interim results. A directional conclusion drawn from
Granite-8B's HumanEval cell alone---that calibration ``actively hurts''---was reversed by
its MBPP cell within the hour. Per-benchmark variance on a single model exceeded the
method difference being characterized.

\subsection{What we would measure next}
\label{sec:next}

Three directions follow directly. \textbf{A second calibrated method} (GPTQ or DWQ) would
test whether ``calibration does not help'' generalizes beyond AWQ, though given the
stability of the head-to-head we would expect it to strengthen rather than change the
conclusion. \textbf{An expanded domain suite} would resolve the study's most intriguing
loose end: Qwen-7B's AWQ arm solves 8 of 26 domain tasks against RTN's 11 and bf16's 12,
while its HumanEval scores differ by 0.6 points---hinting that aggregate pass@1 on
saturated benchmarks may conceal capability loss that harder tasks expose. At $n = 26$
that is an observation, not a result. \textbf{Mixture-of-experts models} are a genuinely
different question: sparse activation and expert routing interact with quantization in
ways dense-model results cannot predict, and the current generation of strong open coding
models is increasingly MoE.

% ── 10. Threats ─────────────────────────────────────────────────────────────
\section{Threats to Validity}
\label{sec:threats}

\begin{itemize}
  \item \textbf{Provenance mismatch on Qwen-3B} (Table~\ref{tab:paired}$^\dagger$).
  \item \textbf{Calibrated arm is partial.} AWQ is measured on 4 of 5 models
    (\S\ref{sec:calibration})---Phi-4-mini is excluded because
    \texttt{mlx\_lm.quant.awq} does not support its architecture. GPTQ and DWQ remain
    unmeasured; the head-to-head is RTN vs AWQ only, so ``calibration does not help'' is
    supported for one calibrated method, not for calibration in general.
  \item \textbf{AWQ and RTN differ in embedding group size} (32 vs 64), so
    \S\ref{sec:calibration}'s head-to-head is not a pure method contrast.
  \item \textbf{$n=1$ greedy.} Justified by the determinism gate
    (\S\ref{sec:determinism}), but a single sample per problem cannot estimate pass@k.
  \item \textbf{Single host.} One machine, one MLX version. No claim about other Apple
    Silicon configurations.
  \item \textbf{Memory figures unusable} (\S\ref{sec:throughput}).
  \item \textbf{Data collection was not first-time-clean.} Three arms of the matrix
    failed twice before completing, for unrelated infrastructure reasons---a runner bug
    that pre-created the conversion output directory (which \texttt{mlx\_lm convert}
    rejects), and a task daemon that silently accepted work without persisting it. All
    three were subsequently re-run and are included in the results above. Neither failure
    affects the validity of any measured number, but both are recorded because a silent
    half-empty matrix is the failure mode this kind of audit is most exposed to.
\end{itemize}

% ── Data availability ───────────────────────────────────────────────────────
\section*{Data Availability}

Per-problem EvalPlus verdict files, the archived result matrices
(\texttt{measured-results-v0313.json} for the unified mlx-lm 0.31.3 runtime and
\texttt{measured-results.json} for the 0.29.1 cross-version replication), the domain task
suite with its reference solutions and 21 negative implementations, and the conversion and
evaluation scripts are available at
\begin{center}
  \url{https://github.com/american-code/ResearchPapers}
\end{center}
This paper's material is under \texttt{papers/mlx-quantization/}, including the two
archived result matrices (\texttt{measured-results-v0313.json} for the unified
mlx-lm 0.31.3 runtime and \texttt{measured-results.json} for the 0.29.1
cross-version replication).
\begin{center}
\end{center}

\bibliographystyle{plainnat}
\bibliography{refs}

\end{document}
